\documentclass[11pt,a4paper]{article}
\usepackage[utf8]{inputenc}
\usepackage[british]{babel}
\usepackage[a4paper,margin=2.3cm]{geometry}
\usepackage{amsmath,amssymb}
\usepackage{graphicx}
\usepackage{tikz}
\usetikzlibrary{calc,angles,quotes,patterns}
\usepackage[most]{tcolorbox}
\usepackage{enumitem}
\usepackage{array}
\usepackage{titlesec}
\usepackage[hidelinks]{hyperref}

\definecolor{bandgreen}{RGB}{46,139,87}
\definecolor{bandblue}{RGB}{30,100,190}
\definecolor{bandamber}{RGB}{220,140,20}
\definecolor{bandpink}{RGB}{210,70,140}
\definecolor{bandpurple}{RGB}{110,50,160}

\newcounter{prob}
% \prob{label}{title}{source}
\newcommand{\prob}[3]{\par\bigskip\refstepcounter{prob}\label{#1}%
  \noindent\textbf{\theprob.\ #2}\hfill{\small\textit{#3}}\par\nopagebreak\smallskip}
\newcommand{\ans}[2]{\par\smallskip\noindent\textbf{\ref{#1}.}\ #2}

\newcommand{\band}[3]{%
  \clearpage
  \begin{tcolorbox}[colback=#1!12,colframe=#1,boxrule=1pt,arc=2pt]
  {\Large\bfseries\color{#1} #2}\hfill{\large #3}
  \end{tcolorbox}}

\newtcolorbox{diarynote}[1]{colback=black!4,colframe=black!45,boxrule=0.6pt,
  arc=2pt,fonttitle=\bfseries\small,title={#1},fontupper=\small}

\setlength{\parindent}{0pt}
\setlength{\parskip}{3pt}
\definecolor{bandgrey}{RGB}{85,85,85}
\begin{document}

\begin{titlepage}
\centering
\vspace*{3cm}
{\Huge\bfseries The Ladies' Diary\par}
\vspace{0.4cm}
{\LARGE After Leybourn, 1816--1825\par}
\vspace{0.8cm}
{\Large Twenty-six problems from ten issues\par}
\vspace{2cm}
{\large\itshape Virtue and Sense with Female Softness join'd,\\
All that subdues and captivates Mankind,\\
In Britain's matchless Fair resplendent shine;\\
They rule Love's Empire by a Right Divine:\\
Justly their Charms th' astonish'd World admires,\\
Whom Royal Charlotte's bright example fires.\par}
\vspace{0.3cm}
{\small The Diary's title page, 1816\par}
\vfill
{\large Working draft\par}
\end{titlepage}

\section*{A word before the problems}

Thomas Leybourn published his collection in 1817, and it stops with the questions answered
in 1816. This booklet takes up where he left off. It is made from a bound volume of ten
issues, for 1816 to 1825, the 113th to the 122nd almanacs, digitised by Google Books. The
new questions in it run from 1299 to 1447. None of them is in Leybourn. Charles Hutton
edited the Diary until 1818, and Olinthus Gregory from 1819.

Each issue printed the answers to the previous year's fifteen questions and set fifteen
new ones. The answers to the questions set in 1825 appeared in 1826, in the next volume, and we
have checked ours against them. Where we could not read a printed answer, the answer
here is our own.

\subsection*{Who was writing in}

Some of the names are famous now. W.~G.~Horner of Bath, whose method of solving equations
is still taught, set three questions. Hamnett Holditch of Lynn, later of Holditch's
theorem, set the cylinders and cones of problem~\ref{w-holditch}. James Thomson,
professor of mathematics in Belfast and the father of Lord Kelvin, set a question on
series. The prize question for 1822 came from M.~Hachette, professor in the Faculty of
Sciences in Paris. A W.~Marrat wrote from New York, and three questions came from a Mr Isaac
Newton, of Lynn and then of Wisbech.

Some are not famous at all. Problem~\ref{w-cone} was answered by Master F.~Wardle of
Macclesfield, who, the editor tells us, ``has been entirely blind from his infancy; but
has for some time evinced a decided turn for mathematical'' studies.

And there were women, more than in any decade since the 1720s.
\begin{itemize}[leftmargin=*,itemsep=1pt]
\item \textbf{Miss Mary Knight} of Boston, almost certainly Wormgate in Boston,
Lincolnshire, set problem~\ref{w-knight} in 1817 and answered it herself. Her solution
to the parabola of problem~\ref{w-parabola} was printed in 1817, and her name is in the
solvers' lists through 1818.
\item \textbf{Miss Susan Jackson} of Mile End, already a solver in Leybourn's last volume,
set problem~\ref{w-jackson} and answered others.
\item \textbf{Miss Sarah Ellis} answered in 1824. In 1825 she appears as \textbf{Mrs Sarah
Clay}, ``formerly Miss S.~Ellis'', answering problems~\ref{w-sphere} and~\ref{w-godward}.
One of the questions she answered was set by a G.~Ellis of Acomb, near York, perhaps a
relation. Marriage, which seems to have ended Nancy Mason's mathematics thirty years
earlier, did not end hers.
\item \textbf{Miss Clara Clay} answered in 1823. In the lists of 1818 there are also a
Miss Mary R---e and a Miss Mary West, perhaps Leybourn's Mary Weston of Gainsborough.
\end{itemize}

\subsection*{How to read the labels}

\textit{Q1314} is the Diary's own question number. \textit{After Q1331} means the question
has been retold; the answer at the back is then for the retelling. Yards, feet and inches
stay, because the answers depend on them. The bands are those of the first book in this
series; there is nothing young enough for Green or Blue. Mechanics goes in the last
section, Beyond A-level.

\subsection*{How this edition was made}
This edition was prepared with the assistance of Claude, an AI model made by Anthropic.
Under the editor's direction, Claude worked from Google Books' machine-read text of the
volume. It drafted the retellings, solutions, notes and introduction, recomputed every
answer, and typeset the text. Selection, editorial decisions and final responsibility for
the text are the editor's.

%================ AMBER ================
\band{bandamber}{Amber}{ages 11--14}

\prob{w-digits}{The square with a two-digit side}{Q1347, by William Burdon}
The side of a square is a two-digit number. Dividing one digit by the other gives the same
as subtracting one from the other. The square of the number times the square of the
number written backwards is 1,016,064. What is the side?

\prob{w-three}{Three numbers}{Q1419, by Henry Ward}
Three numbers add up to 14. Take the differences between the first and second and
between the second and third: these two differences differ by 5. The first times the
third times the cube of the second is 270. Find the numbers.

\prob{w-river}{Rowing against the stream}{Q1435, by Petrarch}
A man rows down a river to a place and back again. Show that it takes him longer when
the river flows than it would on still water, if he rows just as hard.

%================ PINK ================
\band{bandpink}{Pink}{ages 14--16}

\prob{w-knight}{The cone with a slant of 10}{Q1314, by Miss Mary Knight}
The slant side of a cone is 10, and its volume is 94.248. Find its height.

\prob{w-excircle}{The circle outside the triangle}{Q1389, by John Collins}
A circle touches the hypotenuse of a right-angled triangle and the other two sides
extended. Show that its diameter equals the perimeter of the triangle.

\prob{w-sixr}{Three points on a circle}{Q1391, by W.~Rutherford}
Divide a circle into three equal arcs. Join the three points of division to any point on
the circle. Show that the sum of the squares of the three lines is six times the square of
the radius.

\prob{w-arc}{The middle of the arc}{Q1405, by Thomas Baker}
$AB$ is a chord of a circle, $H$ the midpoint of one of its arcs, and $C$ any point of
that arc. Show that $AC\cdot CB + CH^2$ is the same wherever $C$ is.

\prob{w-windows}{The rope across the street}{after Q1331, by William Burdon}
Two windows face each other across a street 17 yards wide, one 3 yards higher than the
other. The shortest rope that can be stretched from one window to the other, touching the
street, is 39 yards long. How high are the windows?

\prob{w-trees}{Two trees in an elliptical garden}{Q1395, by William Burdon}
An elliptical garden has axes of 20 and 15 yards. At its two foci stand trees 8 and 10
yards high. A rope runs from the top of one tree down to the edge of the garden and up to
the top of the other. How short can it be?

\prob{w-cone}{A cone whose volume equals its surface}{Q1404, by Joseph Gritton}
The diameter of the base of a cone equals its slant side, and its volume is numerically
equal to its whole surface, base included. Find its size.

\prob{w-locus}{Dividing lines to a circle}{Q1426, by Isaac Newton of Lynn}
From a fixed point $P$, draw a line $PC$ to any point $C$ of a given circle, and divide it
at $E$ in a fixed ratio. Show that $E$ always lies on a circle.

\prob{w-composite}{Ten numbers, none of them prime}{Q1422, by Investigator}
Find ten consecutive whole numbers, none of which is prime, by a method that would work
for any number of them.

\prob{w-sphere}{The cone and the sphere inside it}{Q1423, by G.~Ellis}
A cone has a volume of $96\pi \approx 301.5936$ cubic inches. Its centre of gravity is
one inch farther from its vertex than the centre of the largest sphere that fits inside
it. Find the cone and the sphere.

%================ PURPLE ================
\band{bandpurple}{Purple}{ages 16--17}

\prob{w-parabola}{The length of a parabola}{Q1299, by George Brown}
The curved length of a parabola, from one end of its base round to the other, is 24, and
its latus rectum is 4. Find its area.

\prob{w-golden}{A triangle in geometric progression}{Q1354, by Charon}
Construct a right-angled triangle whose sides are in geometric progression, and find their
ratio.

\prob{w-log}{A number and its logarithm}{Q1359, by John Baines}
Which number differs least from its common logarithm? Which differs least from its natural
logarithm?

\prob{w-holditch}{Cylinders and cones without end}{Q1338, by Hamnett Holditch}
A cylinder whose diameter equals its height holds a cone with the same base and height.
That cone holds a cylinder of half its height, standing on the same base; that cylinder
holds a cone in the same way, and so on for ever. What do all the cylinders, and all the
cones, hold together?

\prob{w-ladder}{The sliding rod}{Q1401, by W.~Marrat}
A rod of fixed length slides with its top against a vertical line and its foot on a
horizontal line. Find the curve that all its positions touch.

\prob{w-napoleon}{Three equilateral triangles}{Q1439, by W.~Rutherford}
On the three sides of any triangle draw equilateral triangles, all outwards or all
inwards. Show that their centres are the corners of an equilateral triangle.

\prob{w-ellipses}{Two ellipses}{Q1436, by Isaac Newton of Wisbech}
Two ellipses of the same shape have the same centre and their axes in line, one inside the
other. Show that a tangent to the inner one cuts off a segment of the outer one whose area
is always the same.

\prob{w-horner}{A function that undoes itself}{Q1385, by W.~G.~Horner}
For which $a$, $b$, $c$, $d$ does $f(x) = \dfrac{a+bx}{c+dx}$ satisfy $f(f(x)) = x$ for
every $x$?

\prob{w-sixn}{Primes of the form $6n\pm1$}{Q1356, by Henry Atkinson}
Every prime above 3 has the form $6n+1$ or $6n-1$, but not every number of those forms is
prime. Find for which $n$ the number $6n+1$ is composite, and the same for $6n-1$.

\prob{w-grandi}{Divergent series}{Q1442, by Investigator}
By a general method, find the sums of
\[ 1 - 1 + 1 - 1 + \cdots, \qquad 1 + 0 - 1 + 1 + 0 - 1 + \cdots, \qquad
1 + 0 + 0 - 1 + 1 + 0 + 0 - 1 + \cdots. \]

%================ BEYOND ================
\band{bandgrey}{Beyond A-level}{motion}

\prob{w-jackson}{Two balls from a tower}{Q1316, by Miss Susan Jackson}
A ball is dropped from the top of a tower. When it has fallen 30 feet, a second ball is
let go from a point 80 feet below the top. Both reach the ground together. How high is the
tower, and how long was the first ball falling? Take $g = 32\tfrac16$ feet per second per
second.

\prob{w-godward}{The quickest cone}{Q1429, by W.~Godward}
Of all cones with the same volume, find the one down whose slant side a body slides from
the vertex in the least time.

\prob{w-bowl}{The ball in the bowl}{Q1396, by D.~T.~Sheridan}
A smooth ball runs from rest at the rim down the inside of a polished hemispherical bowl.
Where is its downward speed greatest?

%================ ANSWERS ================
\clearpage
\section*{Answers}

\ans{w-digits}{42. Its digits give $4\div2 = 2 = 4 - 2$, and $42\times24 = 1008$, whose
square is 1,016,064. Henry Atkinson pointed out that the last condition cannot tell 42
from 24, since the product is the same either way.}

\ans{w-three}{1, 3 and 10. If the numbers are $x$, $y$, $z$, then $(z-y)-(y-x) = 5$ and
$x+y+z = 14$ give $y = 3$ and $x + z = 11$, and $27xz = 270$ gives $xz = 10$. The editor
printed no solution. Two readers had told him that Mr Ward had sent the same question to
the \textit{Leeds Correspondent}, which had already answered it, and he wrote: ``We wonder
it did not occur to Mr.~Ward that this procedure is inconsistent with the etiquette of
periodical publications.''}

\ans{w-river}{Let the distance be $d$, his speed in still water $v$ and the stream's $w$.
The round trip takes $\dfrac d{v+w} + \dfrac d{v-w} = \dfrac{2dv}{v^2-w^2}$, which is
more than $\dfrac{2d}v$ whenever $w \ne 0$. John Baines and ten others gave this in
1826.}

\ans{w-knight}{About 9.515. The radius is $\sqrt{100 - h^2}$, so
$\tfrac\pi3(100 - h^2)h = 94.248 = 30\pi$, that is $h^3 - 100h + 90 = 0$. Miss Knight
gave 9.51533. Thomas Weatherill noticed that the cubic has a second positive root,
$h \approx 0.9075$: a flat cone with the same slant and volume.}

\ans{w-excircle}{The circle is the excircle opposite the right angle $C$. With centre $O$,
draw $OD$ and $OE$ to the extended legs; $ODCE$ is a square of side $r$. Tangents from $A$
and $B$ give $AD + BE = AB$, so $2r = CD + CE = AC + AD + BC + BE$ is the perimeter. This
was the proof of Master E.~P.~Cooper, a pupil at Wintringham.}

\ans{w-sixr}{Put the circle's centre at the origin and the points of division at angles
$\theta$, $\theta + 120^\circ$, $\theta + 240^\circ$. A point at angle $\phi$ is at squared
distance $2r^2 - 2r^2\cos(\phi - \theta_k)$ from each, and the three cosines add up to
zero. So the sum is $6r^2$. The Rev.~John Sampson noted that for any point at distance $d$
from the centre the sum is $3(r^2 + d^2)$.}

\ans{w-arc}{The constant is $AH^2$. Put $C$ at $A$ to see it. In general, if $HE$ is the
diameter through $H$ and $G$ is the midpoint of $AB$, then $AC\cdot CB + CH^2 = HE\cdot HG
= AH^2$. This was the answer of Mervyn Crawford, a pupil of Mr Herapath.}

\ans{w-windows}{About 19.05 and 16.05 yards. Reflect the lower window in the street. The
shortest rope then runs straight from the upper window to the image, so
$39^2 = 17^2 + (h_1 + h_2)^2$ and $h_1 + h_2 = \sqrt{1232} \approx 35.10$. With
$h_1 - h_2 = 3$, the heights follow. John Davey's printed answer is the same.}

\ans{w-trees}{About 26.91 yards. If the rope reaches the edge at a point $d_1$ from one
focus, it is $d_2 = 20 - d_1$ from the other, since the focal distances of an ellipse add
up to the major axis. So the rope is $\sqrt{d_1^2 + 64} + \sqrt{d_2^2 + 100}$, least when
$d_1 : d_2 = 8 : 10$, and then it is $\sqrt{20^2 + 18^2} = \sqrt{724} \approx 26.907$.
J.~A.~Dotchen gave 26.907248, and remarked that the minor axis is not needed.}

\ans{w-cone}{Base diameter $6\sqrt3 \approx 10.392$, height 9, volume
$81\pi \approx 254.47$. With radius $r$, the slant is $2r$ and the height $r\sqrt3$; then
$\tfrac\pi3r^2\cdot r\sqrt3 = \pi r^2 + 2\pi r^2$ gives $r = 3\sqrt3$. This was Master
Wardle's answer.}

\ans{w-locus}{Let $O$ be the centre and $PE : PC = m : n$. The map that sends each point $X$
to the point dividing $PX$ in that ratio is an enlargement from $P$ with scale factor
$m/n$. It sends the given circle to a circle whose centre divides $PO$ in the same ratio
and whose radius is $m/n$ of the original. This was Anthony Cook's answer.}

\ans{w-composite}{114 to 123. N.~J.~Andrew took the numbers $6n, 6n+1, \dots, 6n+9$; the
only ones that can be prime are $6n+1$, $6n+5$ and $6n+7$, and he made them divisible by
5, 7 and 11, which gives $n = 19$. A method for any length $k$: the numbers
$(k+1)! + 2, \dots, (k+1)! + (k+1)$ are all composite. In fact 113 and 127 are primes with
thirteen composites between them.}

\ans{w-sphere}{There are two cones. The first has base radius 6 and height 8, with an
inscribed sphere of radius 3. The second has base radius about 7.786 and height about
4.751, with a sphere of radius about 2.188. In each, the centre of gravity is $\tfrac34h$
from the vertex and the sphere's centre is $h - \rho$ from it, with
$\rho = rh/(r + \sqrt{r^2+h^2})$. The proposer, G.~Ellis, and Mrs Sarah Clay found the
first. For the second they gave radius 4.123 and height 16.96; see the grey box.}

\ans{w-parabola}{About 80.07. Take the parabola as $y^2 = 4x$. The arc from the vertex to
$(x, y)$ is $\int_0^y\sqrt{1 + t^2/4}\,dt$, and setting this equal to 12 gives
$y \approx 6.2162$ and $x = y^2/4 \approx 9.660$. The area is $\tfrac43xy \approx 80.07$.
J.~A.~Dotchen found 80.0616. Miss Mary Knight used a rule of thumb from Hutton's
\textit{Mensuration}, which gave $x = 9$ and an area of 72. Dotchen pointed out that the
rule fails when the ordinate is much greater than half the latus rectum, as here.}

\ans{w-golden}{The sides go as $1 : \sqrt\varphi : \varphi$, where $\varphi = (1+\sqrt5)/2$,
so the ratio is $\sqrt\varphi \approx 1.272$. If the sides are $a$, $aq$, $aq^2$, Pythagoras
gives $q^4 = q^2 + 1$. The editor printed no solution, because most correspondents had
recognised the question as Problem~39 of the appendix to Simpson's \textit{Algebra}. He
added kindly: ``So many hundred questions have already been published in reference to the
parts of plane triangles, that, unless a person reads and remembers every thing, he can
scarcely avoid similarity to what has gone before.''}

\ans{w-log}{$\log_{10}e \approx 0.4343$, and 1. The difference $x - \log_b x$ is least
where $1 = 1/(x\ln b)$, that is at $x = 1/\ln b$, the modulus of the system. John Collins
and Thomas Martin stated it in just this way.}

\ans{w-holditch}{Each new cylinder has half the height and half the radius of the one
before, so an eighth of its volume. If the first cylinder holds $V$, all the cylinders
hold $V(1 + \tfrac18 + \tfrac1{64} + \cdots) = \tfrac87V$. Each cone is a third of its
cylinder, so the cones hold $\tfrac8{21}V$, and everything together $\tfrac{32}{21}V$.
Thomas Squire's printed answer is the same.}

\ans{w-ladder}{An astroid, $x^{2/3} + y^{2/3} = l^{2/3}$ for a rod of length $l$. The rod's
positions are the lines $x/\cos\theta + y/\sin\theta = l$. Differentiating with respect
to $\theta$ and solving gives $x = l\cos^3\theta$, $y = l\sin^3\theta$. Several
correspondents noted that this was Question 386 in the \textit{Gentleman's Mathematical
Companion}, and the editor printed none of the solutions.}

\ans{w-napoleon}{This is now called Napoleon's theorem. In complex numbers, with
$\omega = e^{2\pi i/3}$, the centre of the equilateral triangle on $BC$ is
$\tfrac13(B + C + A')$, where $A' = -\omega B - \omega^2C$ is its third vertex. A short
calculation shows that the three centres $P$, $Q$, $R$ satisfy
$P + \omega Q + \omega^2R = 0$, which is the condition for an equilateral triangle.
In 1826 Thomas Burn and John Walker gave a proof by congruent triangles.}

\ans{w-ellipses}{Squash the figure along the major axis until both ellipses become
circles. A tangent to the inner circle cuts off a segment of the outer circle whose area
depends only on the two radii. Squashing multiplies every area by the same factor, so the
segments of the ellipses are all equal too. Isaac Brown and James Lugg proved it in 1826
by fluxions.}

\ans{w-horner}{Either $f(x) = x$ itself, or $c = -b$ with $a$ and $d$ arbitrary. Composing
gives $f(f(x)) = \dfrac{a(b+c) + (b^2 + ad)x}{(c^2 + ad) + d(b+c)x}$, which is $x$ for
every $x$ exactly when $(b+c)a = (b+c)d = 0$ and $b^2 = c^2$. Horner's own answer found
all functions whose $n$th repeat is $x$, using a result from Euler.}

\ans{w-sixn}{$6n+1$ is composite exactly when $n = 6pq + p + q$ or $n = 6pq - p - q$ for
whole numbers $p, q \ge 1$, because its factors are both of the form $6k+1$ or both of the
form $6k-1$. $6n-1$ is composite exactly when $n = 6pq + p - q$, from one factor of each
kind. This was Henry Atkinson's own answer.}

\ans{w-grandi}{$\tfrac12$, $\tfrac23$ and $\tfrac34$. Replace each term by the term
times $x^k$ and let $x$ tend to 1. For the third, the series becomes
$(1 - x^3)/(1 - x^4) \to \tfrac34$. In general a series with period $m$, $1, 0, \dots,
0, -1$, sums to $(m-1)/m$. This is the same as averaging the partial sums. Such
``sums'' were argued over from Leibniz to Euler. John Baines junior answered it in 1826 in exactly this way.}

\ans{w-jackson}{About 100.8 feet and 2.50 seconds. After the first ball has fallen 30
feet it moves faster than the second by the constant amount $v_0 = \sqrt{60g}$. It must
make up the 50 feet between them, which takes $50/v_0$ seconds, in which the second ball
falls $\tfrac12g(50/v_0)^2 = 2500/120 = 20\tfrac56$ feet, whatever $g$ is. So the tower is
$100\tfrac56$ feet high. J.~A.~Dotchen gave 100.83 feet and 2.50388 seconds.}

\ans{w-godward}{The height should be $\sqrt2$ times the radius of the base. Sliding from
rest down a slant of length $s$ at height $h$ takes a time whose square is
$2s^2/gh = 2(h^2 + r^2)/gh$. With $r^2h$ fixed, $h + r^2/h$ is least when $h^2 = 2r^2$.
Several solvers, Mrs Sarah Clay among them, gave this; others used the fact that a body
slides down every chord from the top of a vertical circle in the same time.}

\ans{w-bowl}{At $35^\circ16'$ of arc below the rim, where the tangent of the angle is
$1/\sqrt2$. At depth angle $\theta$ the speed is $\sqrt{2gR\sin\theta}$, directed along the
bowl, and its downward part is $\sqrt{2gR\sin\theta}\cos\theta$, which is greatest when
$\tan^2\theta = \tfrac12$. The Diary's answer gives the same arc.}

\begin{diarynote}{When the Diary was wrong}
\textbf{Problem~\ref{w-sphere}.} The second cone printed, radius 4.123 and height 16.96,
satisfies the condition the other way round: its sphere's centre is an inch farther from
the vertex than its centre of gravity. The second cone that fits the question as set has
radius about 7.79 and height about 4.75.

\textbf{Problem~\ref{w-parabola}.} J.~A.~Dotchen's area, 80.0616, is a little low; the true
value is 80.066.
\end{diarynote}

\begin{diarynote}{Where the Diary simply worked differently}
\textbf{Problem~\ref{w-parabola}.} Miss Knight's 72 is not a slip in her working. She used
an approximate rule for the length of a parabola from Hutton's \textit{Mensuration},
the standard book, which is good only for flat parabolas.

\textbf{Problems~\ref{w-golden} and~\ref{w-ladder}.} Both had been published before, in
Simpson's \textit{Algebra} and the \textit{Gentleman's Mathematical Companion}, and the
editor said so.
\end{diarynote}

\vfill
{\small Source: \textit{The Ladies' Diary} for 1816--1825, digitised by Google Books.}

\end{document}
